\chapter{Abdominal Hysterectomy}
\index{Abdominal Hysterectomy}

\section*{Definition and Overview}
An abdominal hysterectomy is an operation to remove the uterus through an incision in the lower abdomen. It is the traditional approach to hysterectomy and is still used when the uterus is large, when keyhole or vaginal routes are not suitable, or when cancer is being treated. The incision is usually a horizontal cut (bikini-line) or occasionally vertical. Abdominal hysterectomy is major surgery with a longer recovery than keyhole or vaginal approaches, but it is safe and highly effective for the conditions it treats.

\section*{Epidemiology}
Abdominal hysterectomy remains a common operation, although the proportion performed by keyhole and vaginal routes has increased substantially. It is most often used for large fibroids, gynaecological cancers, and complex pelvic disease. Around a quarter to a third of local hysterectomies are performed abdominally. Recovery takes 4--6 weeks, longer than for minimally invasive approaches, but the long-term results are equivalent.

\section*{Aetiology and Pathophysiology}
\begin{itemize}
    \item \textbf{Indications:} large fibroids, heavy bleeding, endometriosis, prolapse, and cancers of the uterus, cervix, or ovaries
    \item \textbf{Incision:} a transverse (bikini-line) or vertical lower abdominal incision
    \item \textbf{Extent:} the uterus and cervix are removed; ovaries may also be removed
    \item \textbf{Why abdominal:} large uterine size, cancer surgery, or surgeon preference and safety
    \item \textbf{Pathophysiology:} direct access through the abdominal wall allows the surgeon to see and access the pelvic organs
    \item \textbf{The operation:} the uterus is freed from its attachments and removed, and the incision is closed in layers
\end{itemize}

\section*{Clinical Features}
\begin{itemize}
    \item A horizontal or vertical scar on the lower abdomen
    \item Hospital stay of 2--4 days
    \item Recovery of 4--6 weeks before full activity
    \item Pain managed with medicines and early mobilisation
    \item The end of menstruation and fertility
    \item Surgical menopause if the ovaries are removed
    \item Most symptoms from the underlying condition resolve
\end{itemize}

\section*{Differential Diagnosis}
\begin{itemize}
    \item Keyhole (laparoscopic) hysterectomy, which has faster recovery
    \item Vaginal hysterectomy, which has no abdominal scar
    \item The route is chosen by the indication and patient factors
    \item Alternatives to hysterectomy may be appropriate for benign conditions
    \item The surgeon discusses the best route for each woman
\end{itemize}

\section*{When to Seek Medical Care}
Women should discuss the surgical approach and alternatives with their gynaecologist, and seek urgent care for heavy bleeding, severe pain, or other concerning symptoms.

\textbf{Contact your local emergency services for:}
\begin{itemize}
    \item Heavy vaginal bleeding with faintness
    \item Severe abdominal pain or fever after surgery
    \item Chest pain, breathlessness, or a swollen painful leg (possible blood clot)
    \item Wound problems: spreading redness, discharge, or opening
    \item Difficulty passing urine or opening bowels after surgery
\end{itemize}

\section*{Diagnosis and Investigations}
\begin{itemize}
    \item \textbf{Preoperative assessment:} history, examination, imaging, and fitness for surgery
    \item \textbf{Ultrasound and MRI:} assess uterine size and the underlying condition
    \item \textbf{Biopsy:} for abnormal bleeding or suspected cancer
    \item \textbf{Blood tests and ECG:} preoperative checks
    \item \textbf{Counselling:} discussion of the procedure, recovery, and risks
\end{itemize}

\section*{Management}
\begin{itemize}
    \item \textbf{The operation:} under general anaesthesia, with a lower abdominal incision
    \item \textbf{Pain management:} epidural or patient-controlled analgesia
    \item \textbf{Early mobilisation:} walking from the day of surgery reduces clot risk
    \item \textbf{Anticoagulation:} medicines to prevent blood clots
    \item \textbf{Wound care:} keeping the incision clean and dry
    \item \textbf{Graduated recovery:} increasing activity over 4--6 weeks
    \item \textbf{Follow-up:} review of recovery and pathology results
\end{itemize}

\section*{Complications}
\begin{itemize}
    \item Bleeding and the need for transfusion
    \item Wound infection and slow healing
    \item Damage to the bladder, ureters, or bowel
    \item Blood clots in the legs or lungs
    \item Infection within the pelvis
    \item Hernia at the incision site
    \item Longer recovery than minimally invasive approaches
\end{itemize}

\section*{Prognosis and Outcomes}
Abdominal hysterectomy is highly effective: the symptoms of the underlying condition resolve, and long-term outcomes match other surgical routes. Recovery takes longer than with keyhole or vaginal surgery, but the vast majority of women return to full activity within 6 weeks and experience lasting improvement in quality of life. With modern surgical care, complications are uncommon.

\section*{Prevention and Self-Care}
\begin{itemize}
    \item Discuss all options, including less invasive routes, before surgery
    \item Stop smoking and optimise health before surgery
    \item Arrange help at home for the recovery period
    \item Follow the discharge instructions for wound care
    \item Increase activity gradually as advised
    \item Attend physiotherapy and follow-up
    \item Report fever, wound problems, or leg pain promptly
    \item Seek support for emotional adjustment
\end{itemize}

\section*{Sources and Further Reading}
The following reputable sources provide further information for healthcare students and patients seeking reliable, current advice about abdominal hysterectomy.

\begin{itemize}
    \item Jean Hailes for Women's Health. \textit{Hysterectomy.} https://www.jeanhailes.org.au/
    \item Healthdirect. \textit{Abdominal hysterectomy.} https://www.healthdirect.gov.au/hysterectomy
    \item The Royal Women's Hospital. \textit{Hysterectomy fact sheet.} https://www.thewomens.org.au/
    \item Royal Australian and New Zealand College of Obstetricians and Gynaecologists. \textit{Patient information.} https://ranzcog.edu.au/
    \item NHS. \textit{Hysterectomy: surgery and recovery.} https://www.nhs.uk/conditions/hysterectomy/
    \item Mayo Clinic. \textit{Abdominal hysterectomy.} https://www.mayoclinic.org/
\end{itemize}
